\documentclass[text03.tex]{subfiles}
\begin{document}
\newpage
\section{ボタンを使って LED をつけてみよう}
赤いボタンを\ruby{押}{お}している間 LED が光る HSP プログラムを考えてみましょう。HSP スクリプトエディタで\textasciitilde /03/button\_led.hsp を読み\ruby{込}{こ}み実行してください。センサーボードのボタンのどちらかを押すと LED が 1 つ点灯します。\\

\begin{lstlisting}[caption=\textasciitilde/03/button\_led.hsp,label=button_led.hsp]
#include "hsp3dish.as"		<#blue#;スクリプトの設定を読み込む#>
#include "rpz-gpio.as"		<#blue#;スクリプトの設定を読み込む#>

	gpio 17, 0
	gpio 18, 0
	gpio 22, 0
	gpio 27, 0
   
	redraw 0		<#blue#;画面更新（仮想画面に描画）#>
	font "",30		<#blue#;フォントサイズを決める#>
	pos 20,20		<#blue#;文字の位置を決める#>
	mes "ボタンを押している間 LED が光るよ"		<#blue#;表示する文字を決める#>
	redraw 1		<#blue#;画面更新（実際の画面に描画）#>

*led
	btn1 = gpioin(5) 	<#blue#;ボタンの状態を調べる#>
	if btn1=0 : gpio 18, 1 : else : gpio 18, 0 	<#blue#;ボタンによってLEDを消灯/点灯#>
	wait 10 		<#blue#;0.1 秒待つ#>
	goto *led 		<#blue#;*led に戻る#>
\end{lstlisting}

ここで出てくる新しい命令は \code{gpioin} 命令です。\code{gpioin()}は\code{()}の中にボタンの番号（赤いボタンは
GPIO5 なので５、黒いボタンは GPIO ６なので６）を書きます。\\

\code{gpioin(5)}で赤いボタンの状態（押されているか、押されていないか）を調べることができます。\\
\code{btn1 = gpioin(5)}で、赤いボタンの状態を \code{btn1} という変数に代入しています。赤いボタンは押されている間、変数の数字が 0、押されていない間は 1 になります。\\

赤いボタンが押されている間、LED をつけたい。そこで、もし赤いボタンが押されていたら（変数の
数字が 0 のとき）LED を光らせる、それ以外のときは消すという命令を書きます。もし～は \code{if} 命令
を書きます。

\code{if btn1=0} で \code{btn1} 変数の数字が 0 のとき（赤いボタンが押されているとき）は\code{ }: の次に書かれている命令を実行します。\\
\code{: gpio 18, 1} で GPIO18 の LED を光らせます。\\
\code{: else : gpio 18, 0} では、else（それ以外）のとき、GPIO18 の LED を消しています。\\

\begin{tcolorbox}[title=\useOmetoi]
\begin{enumerate}
\addquiz{ボタンの状態を調べる命令を書きましょう。}
\addex{ターミナルを使って \texttt{button\_led.hsp} のコピーを、\texttt{mybutton\_led.hsp} という名前で作りましょう。}
\addex{黒いボタン(GPIO6)を使って LED を光らせるように、\texttt{mybutton\_led.hsp} のプログラムを変えてみましょう。}
\addex{ 赤いボタン(GPIO5)で LED1(GPIO17)、黒いボタン(GPIO6)で LED3(GPIO22)が光るように、\texttt{mybutton\_led.hsp} のプログラムを変えてみましょう。}
\end{enumerate}
\end{tcolorbox}
\end{document}
